% El informe del evento: para quien organiza, para quien auspicia y para el programa de embajadores
% de Stellar. Lo arma scripts/informe.mjs: escribe datos.tex (los números ya agregados, los textos en
% el idioma pedido y las filas de las tablas) y compila este archivo dos veces con XeLaTeX.
% Cada sección se omite sola si el generador no escribió su marca (\rumboif{seccion.<nombre>}).
\documentclass{rumbo-informe}
\input{datos.tex}
\rumboIdioma
\hypersetup{pdftitle={\V{pdf.titulo}}, pdfauthor={Rumbo}, pdfsubject={\V{pdf.asunto}}, pdfcreator={Rumbo}, pdfproducer={Rumbo}}

% ---------- La portada ----------
% Se dibuja entera en el fondo de la página 1: la marca, el aviso de ejemplo, la brújula, el título
% y el bloque de tinta con las tres cifras.
\renewcommand\rumboPortadaFondo{%
  \node[anchor=west, inner sep=0] at (22,-23) {\includegraphics[height=9mm]{marca/logo-tinta.pdf}};
  \node[anchor=west, inner sep=0, font=\rumboCifra\fontsize{20}{20}\selectfont\color{tinta}] at (33.6,-23) {Rumbo};
  \rumboif{ejemplo}{%
    % En el fondo de la página va un nodo de TikZ y no una tcolorbox: sus transparencias no llegan a la página.
    \node[anchor=north west, fill=amarilloSuave, draw=amarillo, line width=0.6pt, rounded corners=3.2mm, inner xsep=5mm, inner ysep=3.4mm, text width=155mm, align=left] at (22,-34) {%
      {\fontsize{10}{13}\selectfont\bfseries\color{tinta}\V{t.ejemplo.titulo}\par}\vspace{0.8mm}%
      {\fontsize{8.8}{12.5}\selectfont\color{tinta}\V{t.ejemplo.texto}\par}};
  }{}
  \begin{scope}[shift={(158,-100)}]\rumboBrujula{31}{tinta!50!papel}{tinta!85!papel}{tinta}\end{scope}
  \node[anchor=south west, inner sep=0, text width=166mm, align=left] at (22,-198) {%
    {\fontsize{11}{14}\selectfont\bfseries\color{naranjaTexto}\V{t.portada.sobre}\par}\vspace{4mm}%
    {\rumboCifra\fontsize{\V{portada.tam}}{\V{portada.interlinea}}\selectfont\color{tinta}\V{evento.nombre}\par}\vspace{6mm}%
    {\fontsize{11.5}{16}\selectfont\color{tinta}\V{evento.fecha}\par}%
    {\fontsize{11.5}{16}\selectfont\color{tinta2}\V{evento.lugar}\par}%
    {\fontsize{11.5}{16}\selectfont\color{tinta2}\V{t.portada.organiza}\par}};
  \fill[tinta] (0,-212) rectangle (210,-297);
  \V{portada.cifras}
  \draw[line width=0.4pt, color=crema2!35!tinta] (22,-271) -- (188,-271);
  \node[anchor=west, inner sep=0, font=\fontsize{8}{10}\selectfont\color{crema2}] at (22,-279) {\V{t.pie}};
  \node[anchor=east, inner sep=0, font=\fontsize{8}{10}\selectfont\color{crema2}] at (188,-279) {\V{t.portada.corte}};
}
% Una cifra de la portada, sobre la tinta: \rumboCifraPortada{x}{número}{etiqueta}{color del número}
\newcommand\rumboCifraPortada[4]{%
  \node[anchor=south west, inner sep=0, font=\rumboCifra\fontsize{44}{44}\selectfont\color{#4}] at (#1,-246) {#2};
  \node[anchor=north west, inner sep=0, text width=50mm, font=\fontsize{9.5}{12.5}\selectfont\color{crema2}] at (#1,-250.5) {#3};
}

\begin{document}

\null\newpage

% ---------- Resumen en números ----------

\rumboSeccion{\V{t.resumen.titulo}}{\V{t.resumen.bajada}}
\begin{rumboCifras}
\V{resumen.cifras}
\end{rumboCifras}

\rumboif{seccion.barras}{%
  \rumboSubtitulo{\V{t.barras.titulo}}
  \begin{rumboTarjeta}
    \V{barras}
    \vspace{-1mm}\rumboNota{\V{t.barras.nota}}
  \end{rumboTarjeta}
}{}

\rumboif{seccion.llegadas}{%
  \rumboSubtitulo{\V{t.llegadas.titulo}}
  \begin{rumboTarjeta}
    \vspace{1mm}\rumboColumnas{\V{llegadas.n}}{\V{llegadas.columnas}}
    \vspace{2.6mm}\rumboNota{\V{t.llegadas.nota}}
  \end{rumboTarjeta}
}{}

% ---------- Asistencia verificada ----------

\rumboif{seccion.asistencia}{%
  \rumboSeccion{\V{t.asistencia.titulo}}{\V{t.asistencia.bajada}}
  \begin{rumboCifras}[2]
    \begin{tcolorbox}[rumbo]
      {\fontsize{9.2}{12}\selectfont\bfseries\V{t.asistencia.verificados}\par}\vspace{2mm}
      {\rumboCifra\fontsize{40}{41}\selectfont\color{tinta}\V{asistencia.verificados}\par}\vspace{2.4mm}
      \V{asistencia.chip}\par\vspace{2mm}
      {\fontsize{8}{11}\selectfont\color{apagado}\V{t.asistencia.minimo}\par}
    \end{tcolorbox}
    \begin{tcolorbox}[rumbo]
      {\fontsize{9.2}{12}\selectfont\bfseries\V{t.minutos.titulo}\par}\vspace{2mm}
      \rumboif{minutos.valor}{%
        {\rumboCifra\fontsize{40}{41}\selectfont\color{tinta}\V{minutos.valor}\par}\vspace{2.4mm}
        \V{minutos.chip}\par\vspace{2mm}
      }{%
        \begin{tikzpicture}[x=1mm, y=1mm]
          \draw[line width=0.6pt, color=tinta2, rounded corners=1.6mm] (0,0) rectangle (34,13);
          \node[anchor=west, inner sep=0, font=\fontsize{10}{12}\selectfont\color{apagado}] at (38,4.5) {\V{t.minutos.unidad}};
        \end{tikzpicture}\par\vspace{2.6mm}
      }%
      {\fontsize{8}{11}\selectfont\color{apagado}\V{t.minutos.nota}\par}
    \end{tcolorbox}
  \end{rumboCifras}
  \rumboif{asistencia.pendientes}{\par\vspace{3mm}\rumboNota{\V{asistencia.pendientes}}}{}

  \rumboSubtitulo{\V{t.asistencia.como.titulo}}
  \rumboParrafo{\V{t.asistencia.como}}
  \rumboif{contratos}{%
    \vspace{4mm}
    \begin{rumboTarjeta}
      \V{contratos}
      \rumboif{asistencia.contraste}{\par\vspace{3mm}\rumboNota{\V{asistencia.contraste}}}{}
    \end{rumboTarjeta}
  }{}

  \rumboif{seccion.transacciones}{%
    \rumboSubtitulo{\V{t.tx.titulo}}
    \begin{rumboTarjeta}
      \V{transacciones}
      \vspace{3.4mm}\rumboNota{\V{t.tx.nota}}
    \end{rumboTarjeta}
  }{}
}{}

% ---------- La comida y lo que se entregó ----------

\rumboif{seccion.comida}{%
  \rumboSeccion{\V{t.comida.titulo}}{\V{t.comida.bajada}}
  \begin{rumboTarjeta}
    \V{comida.tabla}
  \end{rumboTarjeta}
  \rumboif{seccion.gasto}{%
    \vspace{4mm}
    \begin{rumboTarjeta}
      \leavevmode\begin{tabularx}{\linewidth}{@{}X@{\hspace{8mm}}X@{\hspace{8mm}}X@{}}
        {\fontsize{8.6}{11}\selectfont\bfseries\V{t.gasto.total}} &
        {\fontsize{8.6}{11}\selectfont\bfseries\V{t.gasto.persona}} &
        {\fontsize{8.6}{11}\selectfont\bfseries\V{t.gasto.tope}} \\[1.2mm]
        {\rumboCifra\fontsize{26}{28}\selectfont\V{gasto.total}} &
        {\rumboCifra\fontsize{26}{28}\selectfont\V{gasto.persona}} &
        {\rumboCifra\fontsize{26}{28}\selectfont\V{gasto.tope}} \\
      \end{tabularx}
      \par\vspace{7mm}
      \rumboMedida{\V{gasto.fraccion}}{\V{gasto.fracciontope}}{\V{gasto.color}}{\V{t.gasto.marca}}
      \vspace{4mm}\V{gasto.chip}\par\vspace{2mm}
      \rumboNota{\V{t.gasto.nota}}
    \end{rumboTarjeta}
  }{}
}{}

% ---------- Para quien auspicia ----------

\rumboif{seccion.auspicio}{%
  \rumboSeccion{\V{t.auspicio.titulo}}{\V{t.auspicio.bajada}}
  \begin{rumboCifras}
  \V{auspicio.cifras}
  \end{rumboCifras}
  \vspace{3mm}
  \rumboNota{\V{t.auspicio.nota}}
}{}

% ---------- El sorteo ----------

\rumboif{seccion.sorteo}{%
  \rumboSeccion{\V{t.sorteo.titulo}}{\V{t.sorteo.bajada}}
  \begin{rumboCifras}[2]
  \V{sorteo.cifras}
  \end{rumboCifras}
  \rumboif{sorteo.premios}{\vspace{4mm}\begin{rumboTarjeta}\V{sorteo.premios}\end{rumboTarjeta}}{}
  \rumboif{sorteo.comprobante}{%
    \rumboSubtitulo{\V{t.sorteo.como.titulo}}
    \rumboParrafo{\V{t.sorteo.como}}
    \vspace{2.4mm}
    \V{sorteo.comprobante}
  }{}
}{}

% ---------- Cómo verificar este informe ----------

\rumboSeccion{\V{t.verificar.titulo}}{}
\begin{tcolorbox}[rumboOscura]
  {\fontsize{10}{15.5}\selectfont\V{t.verificar.texto}\par}
  \vspace{3mm}
  {\fontsize{8.4}{12}\selectfont\color{crema2}\V{t.verificar.corte}\par}
\end{tcolorbox}

\end{document}
